% ============================================================
% DISPERSION CLASS 1
% Statistics - Module 3
% Topics: Range, Quartiles, IQR, Quartile Deviation
% Individual, Discrete and Continuous Series
% ============================================================

\documentclass[12pt,a4paper]{article}

\usepackage[margin=1in]{geometry}
\usepackage{amsmath,amssymb}
\usepackage{array}
\usepackage{booktabs}
\usepackage{enumitem}
\usepackage{fancyhdr}
\usepackage[T1]{fontenc}
\usepackage{lmodern}
\usepackage{xcolor}

\pagestyle{fancy}
\fancyhf{}
\lhead{Statistics -- Module 3}
\rhead{Dispersion Class 1}
\cfoot{\thepage}

\setlength{\parindent}{0pt}
\setlength{\parskip}{6pt}

\newcommand{\boxformula}[1]{%
\[
\boxed{\displaystyle #1}
\]
}

\title{\textbf{Statistics -- Module 3}\\
\Large Measures of Dispersion -- Class 1}
\author{Exam Preparation Notes}
\date{}

\begin{document}

\maketitle

\section*{What is Dispersion?}

\textbf{Dispersion} tells us how much the observations in a dataset are
spread out or scattered.

The topics covered in this class are:

\begin{enumerate}
    \item Range
    \item Coefficient of Range
    \item Quartiles
    \item Interquartile Range (IQR)
    \item Quartile Deviation (QD)
    \item Coefficient of Quartile Deviation
    \item Quartiles in Individual Series
    \item Quartiles in Discrete Series
    \item Quartiles in Continuous Series
\end{enumerate}

% ============================================================
\section{1. Range}

\subsection*{Meaning}

Range is the difference between the largest and smallest observations.

\boxformula{\text{Range}=L-S}

where:
\begin{itemize}
    \item $L$ = largest value
    \item $S$ = smallest value
\end{itemize}

\subsection*{Worked Question 1}

\textbf{Find the range:}
\[
12,\;7,\;19,\;3,\;15,\;8
\]

Largest value:
\[
L=19
\]

Smallest value:
\[
S=3
\]

Therefore:
\[
R=19-3=\boxed{16}
\]

\subsection*{Worked Question 2: Negative Values}

\textbf{Find the range:}
\[
-8,\;-3,\;4,\;11,\;-15,\;7
\]

Largest value:
\[
L=11
\]

Smallest value:
\[
S=-15
\]

Therefore:
\[
R=11-(-15)
\]
\[
=11+15
\]
\[
\boxed{R=26}
\]

\subsection*{Important}

The data does not have to be sorted. Simply identify the maximum and minimum values.

\subsection{Coefficient of Range}

The coefficient of range expresses the range relative to the extreme values.

\boxformula{
\text{Coefficient of Range}
=
\frac{L-S}{L+S}
}

\subsection*{Worked Question 3}

\textbf{For the data}
\[
10,\;20,\;30,\;40,\;50
\]
\textbf{find the range and coefficient of range.}

Largest:
\[
L=50
\]

Smallest:
\[
S=10
\]

Range:
\[
R=50-10=\boxed{40}
\]

Coefficient of Range:
\[
\frac{50-10}{50+10}
=
\frac{40}{60}
=
\frac{2}{3}
\]

Therefore:
\[
\boxed{\text{Coefficient of Range}\approx0.667}
\]

% ============================================================
\section{2. Quartiles}

Quartiles divide ordered data into four parts.

\[
Q_1=25^{\text{th}}\text{ percentile}
\]

\[
Q_2=50^{\text{th}}\text{ percentile}=\text{Median}
\]

\[
Q_3=75^{\text{th}}\text{ percentile}
\]

Conceptually:

\[
\text{Minimum}
\quad|\quad
Q_1
\quad|\quad
Q_2
\quad|\quad
Q_3
\quad|\quad
\text{Maximum}
\]

The four sections each contain approximately $25\%$ of the observations.

\subsection*{Important Convention Used in Class}

For individual/ungrouped data, the class uses:

\boxformula{
Q_1=\text{size of }\frac{N+1}{4}\text{th item}
}

\boxformula{
Q_3=\text{size of }\frac{3(N+1)}{4}\text{th item}
}

Here $N$ is the number of observations.

% ============================================================
\section{3. Interquartile Range}

The Interquartile Range measures the spread of the middle $50\%$ of the data.

\boxformula{IQR=Q_3-Q_1}

It is the distance from $Q_1$ to $Q_3$.

\[
Q_1
\quad\underbrace{\rule{3cm}{0.4pt}}_{\text{middle 50\%}}
\quad Q_3
\]

Unlike ordinary range, IQR is less affected by extreme values.

% ============================================================
\section{4. Quartile Deviation}

Quartile Deviation is also called the \textbf{Semi-Interquartile Range}.

It is half of the Interquartile Range.

\boxformula{
QD=\frac{Q_3-Q_1}{2}
}

Since:
\[
IQR=Q_3-Q_1
\]

we can also write:
\[
\boxed{QD=\frac{IQR}{2}}
\]

% ============================================================
\section{5. Coefficient of Quartile Deviation}

The coefficient of quartile deviation is a relative measure of quartile dispersion.

\boxformula{
\text{Coefficient of QD}
=
\frac{Q_3-Q_1}{Q_3+Q_1}
}

% ============================================================
\section{6. Individual Series -- Worked Example}

\textbf{Question:} From the following data, compute Interquartile Range,
Quartile Deviation and Coefficient of Quartile Deviation.

\[
24,\;7,\;11,\;9,\;17,\;3,\;20,\;14,\;4,\;22,\;27
\]

\subsection*{Step 1: Arrange in ascending order}

\[
3,\;4,\;7,\;9,\;11,\;14,\;17,\;20,\;22,\;24,\;27
\]

There are:
\[
N=11
\]

Number the observations:

\[
\begin{array}{c|ccccccccccc}
\text{Position}
&1&2&3&4&5&6&7&8&9&10&11\\
\hline
X
&3&4&7&9&11&14&17&20&22&24&27
\end{array}
\]

\subsection*{Step 2: Find $Q_1$}

\[
Q_1=\frac{N+1}{4}\text{th item}
\]

\[
=\frac{11+1}{4}
=\frac{12}{4}
=3^{\text{rd}}\text{ item}
\]

The 3rd item is $7$.

\[
\boxed{Q_1=7}
\]

\subsection*{Step 3: Find $Q_3$}

\[
Q_3=\frac{3(N+1)}{4}\text{th item}
\]

\[
=\frac{3(11+1)}{4}
=\frac{36}{4}
=9^{\text{th}}\text{ item}
\]

The 9th item is $22$.

\[
\boxed{Q_3=22}
\]

\subsection*{Step 4: Interquartile Range}

\[
IQR=Q_3-Q_1
\]

\[
=22-7
\]

\[
\boxed{IQR=15}
\]

\subsection*{Step 5: Quartile Deviation}

\[
QD=\frac{Q_3-Q_1}{2}
\]

\[
=\frac{22-7}{2}
=\frac{15}{2}
\]

\[
\boxed{QD=7.5}
\]

\subsection*{Step 6: Coefficient of QD}

\[
\text{Coefficient of QD}
=
\frac{Q_3-Q_1}{Q_3+Q_1}
\]

\[
=
\frac{22-7}{22+7}
=
\frac{15}{29}
\]

\[
\boxed{\text{Coefficient of QD}\approx0.52}
\]

% ============================================================
\section{7. Discrete Series}

\subsection*{What is Discrete Data?}

In a discrete frequency distribution, we have:

\begin{itemize}
    \item $X$ = value
    \item $f$ = frequency (how many times the value occurs)
    \item $cf$ = cumulative frequency
\end{itemize}

Example:

\[
\begin{array}{c|c|c}
X&f&cf\\
\hline
10&7&7\\
20&18&25\\
30&15&40\\
40&25&65\\
50&30&95\\
60&20&115\\
70&16&131\\
80&7&138\\
90&2&140
\end{array}
\]

\subsection*{Meaning of Frequency}

If:
\[
X=10,\quad f=7
\]

it means that the value $10$ occurs $7$ times.

\subsection*{Meaning of Cumulative Frequency}

\textbf{cf} is the running total of frequencies.

For example:
\[
7+18=25
\]

\[
25+15=40
\]

Therefore, the first three cumulative frequencies are:
\[
7,\;25,\;40
\]

\subsection*{Why is cf useful?}

The cumulative frequency tells us where each observation lies in the ordered dataset.

For the above table:

\[
\begin{array}{c|c}
\text{X}&\text{Observation positions}\\
\hline
10&1\text{--}7\\
20&8\text{--}25\\
30&26\text{--}40\\
40&41\text{--}65\\
50&66\text{--}95\\
60&96\text{--}115\\
70&116\text{--}131\\
80&132\text{--}138\\
90&139\text{--}140
\end{array}
\]

\textbf{Key idea:} The class/value tells us \emph{what the observations are}; the
frequency tells us \emph{how many}; cumulative frequency tells us
\emph{how many observations have been covered so far}.

% ============================================================
\subsection{Worked Example: Discrete Series}

\textbf{Question:} Calculate Quartile Deviation and its coefficient.

\[
\begin{array}{c|ccccccccc}
X&10&20&30&40&50&60&70&80&90\\
\hline
f&7&18&15&25&30&20&16&7&2
\end{array}
\]

\subsection*{Step 1: Find $N$}

\[
N=\sum f
\]

\[
=7+18+15+25+30+20+16+7+2
\]

\[
\boxed{N=140}
\]

\subsection*{Step 2: Construct cf}

\[
\begin{array}{c|c|c}
X&f&cf\\
\hline
10&7&7\\
20&18&25\\
30&15&40\\
40&25&65\\
50&30&95\\
60&20&115\\
70&16&131\\
80&7&138\\
90&2&140
\end{array}
\]

\subsection*{Step 3: Find $Q_1$ position}

Using the class method:

\[
Q_1=\frac{N+1}{4}\text{th item}
\]

\[
=\frac{140+1}{4}
=\frac{141}{4}
=35.25^{\text{th}}\text{ item}
\]

Now find the first cf that is equal to or greater than $35.25$.

\[
7,\;25,\;\boxed{40}
\]

The corresponding value of $X$ is $30$.

Therefore:
\[
\boxed{Q_1=30}
\]

Interpretation: the $35.25^{\text{th}}$ observation has the value $30$.

\subsection*{Step 4: Find $Q_3$ position}

\[
Q_3=\frac{3(N+1)}{4}\text{th item}
\]

\[
=\frac{3(141)}{4}
=\frac{423}{4}
=105.75^{\text{th}}\text{ item}
\]

Find the first cf that is equal to or greater than $105.75$:

\[
7,\;25,\;40,\;65,\;95,\;\boxed{115}
\]

The corresponding value is $60$.

Therefore:
\[
\boxed{Q_3=60}
\]

\subsection*{Step 5: Quartile Deviation}

\[
QD=\frac{Q_3-Q_1}{2}
\]

\[
=\frac{60-30}{2}
\]

\[
\boxed{QD=15}
\]

\subsection*{Step 6: Coefficient of QD}

\[
\text{Coefficient of QD}
=
\frac{Q_3-Q_1}{Q_3+Q_1}
\]

\[
=
\frac{60-30}{60+30}
=
\frac{30}{90}
\]

\[
\boxed{\text{Coefficient of QD}=0.33}
\]

\subsection*{Important Discrete-Series Trick}

For discrete data:

\[
\boxed{
\text{Position}\rightarrow cf\rightarrow X
}
\]

That is:

\begin{enumerate}
    \item Calculate the quartile position.
    \item Look down the cf column.
    \item Find the first cf that reaches or passes the position.
    \item Take the corresponding $X$ value.
\end{enumerate}

% ============================================================
\section{8. Continuous Series}

\subsection*{What is Continuous Data?}

In a continuous frequency distribution, values are grouped into class intervals.

Example:

\[
\begin{array}{c|c}
\text{Class interval}&f\\
\hline
5\text{--}7&14\\
8\text{--}10&24\\
11\text{--}13&38\\
14\text{--}16&20\\
17\text{--}19&4
\end{array}
\]

Here, $f$ tells us how many observations belong to each class.

For example:

\[
5\text{--}7,\quad f=14
\]

means there are $14$ observations whose values lie in the $5$--$7$ class.

It does \textbf{not} mean that observations numbered 5 through 7 are being referred to.

% ============================================================
\subsection{Worked Example: Continuous Series}

\textbf{Question:} Calculate Quartile Deviation and its coefficient from:

\[
\begin{array}{c|ccccc}
\text{Class interval}
&5\text{--}7&8\text{--}10&11\text{--}13&14\text{--}16&17\text{--}19\\
\hline
f&14&24&38&20&4
\end{array}
\]

\subsection*{Step 1: Find $N$}

\[
N=14+24+38+20+4
\]

\[
\boxed{N=100}
\]

\subsection*{Step 2: Construct cf}

\[
\begin{array}{c|c|c}
\text{Class}&f&cf\\
\hline
5\text{--}7&14&14\\
8\text{--}10&24&38\\
11\text{--}13&38&76\\
14\text{--}16&20&96\\
17\text{--}19&4&100
\end{array}
\]

\subsection*{Step 3: Find the $Q_1$ class}

For continuous series:

\[
\frac N4=\frac{100}{4}=25
\]

We need the 25th observation.

Cumulative frequency:

\[
\begin{array}{c|c}
\text{Class}&cf\\
\hline
5\text{--}7&14\\
8\text{--}10&38\\
11\text{--}13&76
\end{array}
\]

The 25th observation falls between cf $14$ and cf $38$.

Therefore:

\[
\boxed{Q_1\text{ class}=8\text{--}10}
\]

\subsection*{Step 4: Continuous $Q_1$ Formula}

\boxformula{
Q_1
=
L+
\left(
\frac{\frac N4-cf}{f}
\right)i
}

where:
\begin{itemize}
    \item $L$ = lower boundary of the $Q_1$ class
    \item $cf$ = cumulative frequency before the $Q_1$ class
    \item $f$ = frequency of the $Q_1$ class
    \item $i$ = class width
\end{itemize}

For the class $8$--$10$:

\[
L=7.5,\qquad cf=14,\qquad f=24,\qquad i=3
\]

Therefore:

\[
Q_1
=
7.5+
\left(
\frac{25-14}{24}
\right)3
\]

\[
=
7.5+\frac{11}{24}\times3
\]

\[
=
7.5+1.375
\]

\[
\boxed{Q_1=8.875}
\]

\textbf{Note:} The class slide used in class appears to show $8.75$, but
$7.5+\frac{11}{24}\times3$ gives $8.875$. The displayed answer appears to
contain an arithmetic/calculation inconsistency.

% ============================================================
\subsection*{Step 5: Find the $Q_3$ class}

\[
\frac{3N}{4}
=
\frac{3(100)}4
=
75
\]

We need the 75th observation.

Cumulative frequencies:

\[
38<75\leq76
\]

Therefore:

\[
\boxed{Q_3\text{ class}=11\text{--}13}
\]

\subsection*{Step 6: Continuous $Q_3$ Formula}

\boxformula{
Q_3
=
L+
\left(
\frac{\frac{3N}{4}-cf}{f}
\right)i
}

For the class $11$--$13$:

\[
L=10.5,\qquad cf=38,\qquad f=38,\qquad i=3
\]

Therefore:

\[
Q_3
=
10.5+
\left(
\frac{75-38}{38}
\right)3
\]

\[
=
10.5+\frac{37}{38}\times3
\]

\[
\approx10.5+2.921
\]

\[
\boxed{Q_3\approx13.42}
\]

% ============================================================
\subsection{Understanding Class Boundaries}

This is important for continuous series.

The classes are:

\[
5\text{--}7,\quad
8\text{--}10,\quad
11\text{--}13,\quad
14\text{--}16,\quad
17\text{--}19
\]

There is a gap of $1$ between consecutive class limits.

To obtain continuous boundaries, place the boundary halfway between them.

Between $10$ and $11$:

\[
\frac{10+11}{2}=10.5
\]

Therefore:

\[
11\text{--}13
\quad\longrightarrow\quad
10.5\text{--}13.5
\]

Hence the lower boundary of $11$--$13$ is:

\[
\boxed{L=10.5}
\]

Similarly:

\[
8\text{--}10
\quad\longrightarrow\quad
7.5\text{--}10.5
\]

Thus:

\[
\boxed{\text{Lower boundary}=\text{lower class limit}-0.5}
\]

for these integer-width classes.

% ============================================================
\subsection*{Step 7: Quartile Deviation}

Using the calculated values:

\[
Q_1=8.875,\qquad Q_3\approx13.42
\]

\[
QD=\frac{Q_3-Q_1}{2}
\]

\[
=
\frac{13.42-8.875}{2}
\]

\[
\approx\boxed{2.27}
\]

\subsection*{Step 8: Coefficient of QD}

\[
\text{Coefficient of QD}
=
\frac{Q_3-Q_1}{Q_3+Q_1}
\]

\[
=
\frac{13.42-8.875}{13.42+8.875}
\]

\[
\approx\boxed{0.20}
\]

If following the displayed classroom value $Q_1=8.75$, the slide obtains
approximately $QD=2.34$ and coefficient $0.21$. The discrepancy originates
from the displayed $Q_1$ calculation.

% ============================================================
\section{9. Master Formula Sheet}

\subsection*{Range}

\boxformula{R=L-S}

\subsection*{Coefficient of Range}

\boxformula{
\text{Coefficient of Range}
=
\frac{L-S}{L+S}
}

\subsection*{Individual Series: $Q_1$}

\boxformula{
Q_1=\text{size of }\frac{N+1}{4}\text{th item}
}

\subsection*{Individual Series: $Q_3$}

\boxformula{
Q_3=\text{size of }\frac{3(N+1)}{4}\text{th item}
}

\subsection*{Interquartile Range}

\boxformula{IQR=Q_3-Q_1}

\subsection*{Quartile Deviation}

\boxformula{
QD=\frac{Q_3-Q_1}{2}
}

\subsection*{Coefficient of Quartile Deviation}

\boxformula{
\text{Coefficient of QD}
=
\frac{Q_3-Q_1}{Q_3+Q_1}
}

\subsection*{Continuous Series: $Q_1$}

\boxformula{
Q_1
=
L+
\left(
\frac{\frac N4-cf}{f}
\right)i
}

\subsection*{Continuous Series: $Q_3$}

\boxformula{
Q_3
=
L+
\left(
\frac{\frac{3N}{4}-cf}{f}
\right)i
}

% ============================================================
\section{10. Exam Decision Guide}

When you see a question, first identify the type of data.

\begin{center}
\begin{tabular}{p{4cm}p{8cm}}
\toprule
\textbf{Type} & \textbf{What to do} \\
\midrule
Individual data &
Arrange values in ascending order and use the $(N+1)$ positional formulas.\\
\addlinespace
Discrete frequency data &
Calculate $N$, construct cf, find the quartile position, then locate the
corresponding $X$ using cf.\\
\addlinespace
Continuous frequency data &
Calculate $N$, construct cf, identify the quartile class, then use the
continuous-series quartile formula.\\
\bottomrule
\end{tabular}
\end{center}

\subsection*{Discrete Data Shortcut}

\[
\boxed{\text{Position}\rightarrow cf\rightarrow X}
\]

\subsection*{Continuous Data Shortcut}

\[
\boxed{\text{Position}\rightarrow\text{Quartile class}
\rightarrow L,cf,f,i\rightarrow Q_1/Q_3}
\]

% ============================================================
\section{11. Final Memory Sheet}

\begin{enumerate}[label=\textbf{\arabic*.}]
    \item \textbf{Range} = maximum minus minimum.
    \item \textbf{Coefficient of Range} = $(L-S)/(L+S)$.
    \item $Q_1$ represents the 25th percentile.
    \item $Q_2$ is the median.
    \item $Q_3$ represents the 75th percentile.
    \item \textbf{IQR} = $Q_3-Q_1$.
    \item \textbf{QD} = $(Q_3-Q_1)/2$.
    \item \textbf{Coefficient of QD} = $(Q_3-Q_1)/(Q_3+Q_1)$.
    \item In discrete data, \textbf{cf tells you where an observation falls}.
    \item In continuous data, use the \textbf{quartile class} and the $L,cf,f,i$
    formula.
    \item For the class intervals used here, the lower boundary is the lower
    limit minus $0.5$.
\end{enumerate}

\section*{12. Exam Workflow}

For any quartile-dispersion problem:

\[
\boxed{
\text{Identify data type}
\rightarrow
\text{Find }N
\rightarrow
\text{Find }Q_1,Q_3
\rightarrow
IQR
\rightarrow
QD
\rightarrow
\text{Coefficient}
}
\]

\textbf{Do not panic when you see a frequency table.}
The mathematics is mostly addition, locating a position, subtraction,
division, and substitution into a formula.

\end{document}
